% !TeX root = ../slides.tex
% 04-environment.tex — implementation details of the custom Gymnasium env.

\section{Environment Implementation}

\begin{frame}{Custom Gymnasium Environment}
  \texttt{MinesweeperEnv(gym.Env)} wraps a pure game-logic class
  \texttt{Game}, decoupled from the Gymnasium API and from rendering:
  \begin{itemize}
    \item \texttt{Game}: mine placement, adjacency counts, iterative
      flood-fill uncovering of empty regions (stack-based, avoids deep
      recursion), win/loss detection.
    \item \texttt{MinesweeperEnv}: \texttt{reset}/\texttt{step}/
      \texttt{render}/\texttt{close}, reward assignment, episode
      statistics.
    \item Optional \texttt{human} render mode via Pygame, fully
      independent from the agent's observation.
  \end{itemize}
\end{frame}

\begin{frame}{Design Choices for Stable Training}
  \begin{itemize}
    \item \textbf{Safe first move}: if the first action hits a mine, the
      board is regenerated from the same RNG (not reseeded) until the
      first move is safe — avoids penalizing an unavoidable loss and
      keeps the training signal stable.
    \item \textbf{Mines never revealed to the agent}: even at a losing
      terminal step, the mine stays $-1$ only in the internal board; the
      \texttt{player\_board} keeps it as $-2$. Loss is communicated
      through the terminal reward alone, keeping the observation
      alphabet at 10 symbols (consistent with the planned one-hot
      encoding).
    \item \textbf{Guess detection}: a safe move is a guess if \emph{all}
      its neighbors are still unrevealed — used only for reward shaping,
      not for masking actions.
  \end{itemize}
\end{frame}
